\section{Directional Capability Loss and BRF}
\label{sec:theorical}

% The failure mode we study is not ordinary underperformance on a difficult reverse
% task. A hard inverse can arise because the reverse mapping is underspecified,
% structurally complex, or poorly represented in the training data. Directional
% capability loss is a stronger condition. It arises when forward fine-tuning
% causes the model to generate in the trained direction even when the prompt
% requests the inverse. The distinction has practical weight in software
% engineering, where forward and reverse tasks are operationally adjacent. A model
% that is competent only in the trained direction can still be unsafe inside a
% workflow that depends on rollback, inspection, or reverse analysis.


This section defines the phenomenon we study and the framework we use to measure it.
We first separate directional capability loss from ordinary difficulty on a reverse task, and then introduce BRF.

A model can perform poorly on a reverse task for two different reasons.
The first is that the inverse is simply hard: the reverse mapping may be underspecified, structurally complex, or rare in the training data.
In that case, the model fails because the task exceeds its ability, before or after fine-tuning.
The second reason is the one we study.
Forward fine-tuning can cause a model to keep generating in the trained direction even when the prompt explicitly requests the inverse.
We call this directional capability loss.
It is a stronger condition than difficulty, because it concerns what fine-tuning changes rather than what the task demands.
Section~\ref{sec:bg-motivation} discusses why this distinction matters in software engineering workflows.

Let $T$ be a semantics-preserving program transformation that maps an original
program $C_{\text{orig}}$ to a transformed program $C_{\text{trans}}$, and let
$M_T$ be a model fine-tuned only on the forward direction of $T$. BRF evaluates
$M_T$ on the forward mapping $C_{\text{orig}} \rightarrow C_{\text{trans}}$ and on
the reverse mapping $C_{\text{trans}} \rightarrow C_{\text{pred}}$. Reverse
success cannot be defined by execution alone. Since $C_{\text{trans}}$ is already
semantically equivalent to $C_{\text{orig}}$, a model that copies the transformed
input back will still compile and pass the tests. Reverse evaluation must
therefore ask whether the output moves toward a readable program equivalent to
the original (obfuscation case), rather than remaining at the transformed endpoint.


We measure that asymmetry with a bidirectional gap. For a given model and
transformation, BRF computes forward accuracy
$\mathrm{Acc}_{\text{fwd}}(M_T)$ and reverse recovery
$\mathrm{Acc}_{\text{rev}}(M_T)$ and defines:

\begin{equation}
\Delta_{\text{BRF}}(M_T) = \mathrm{Acc}_{\text{fwd}}(M_T) - \mathrm{Acc}_{\text{rev}}(M_T).
\end{equation}

A large gap does not, on its own, diagnose the phenomenon, because some inverse
tasks are harder than their forward counterparts. We classify the pattern as the
\textit{inversion curse} only when three conditions hold jointly: \bcircle{1} Reverse accuracy is near
zero, 
\bcircle{2} the outputs remain measurably closer to the transformed input than to the
original program, and 
\bcircle{3} stronger prompting does not improve the reverse result.
Under these conditions the model behaves as though adaptation has bound its
generation policy to the trained direction.

BRF is deliberately lightweight. It needs no second model, no architectural
change, and no reverse-direction training data. It adds a single evaluation pass
and surfaces a capability that forward validation leaves untested. The remainder
of the paper instantiates BRF on obfuscation and deobfuscation. The
diagnostic applies to any paired software engineering task that satisfies
Definition~\ref{def:pair}, including translation and back-translation and
compilation and decompilation.
